\chapter{Statements}
\label{ch:statements}

Statements perform actions. Unlike expressions, statements do not produce values.

\section{Variable Declarations}

\subsection{With Inference}

\begin{lstlisting}
x = 42                  // Type inferred as int
name = "Haira"          // Type inferred as string
items = [1, 2, 3]       // Type inferred as []int
\end{lstlisting}

\subsection{With Explicit Type}

\begin{lstlisting}
x: int = 42
name: string = "Haira"
count: i32 = 100
\end{lstlisting}

\subsection{Multiple Declarations}

\begin{lstlisting}
a, b = 1, 2
x, y, z = compute()     // Destructuring tuple
\end{lstlisting}

\subsection{Uninitialized Variables}

\begin{lstlisting}
x: int                  // Must be assigned before use
\end{lstlisting}

\begin{warning}
Using an uninitialized variable is a compile-time error.
\end{warning}

\section{Assignment Statements}

\begin{lstlisting}
x = 10                  // Simple assignment
x += 5                  // Add and assign (x = x + 5)
x -= 3                  // Subtract and assign
x *= 2                  // Multiply and assign
x /= 4                  // Divide and assign
\end{lstlisting}

\subsection{Multiple Assignment}

\begin{lstlisting}
a, b = b, a             // Swap values
x, y = get_coordinates()
quotient, remainder = divide(17, 5)
\end{lstlisting}

\subsection{Field Assignment}

\begin{lstlisting}
user.name = "Bob"
point.x = 10.0
\end{lstlisting}

\subsection{Index Assignment}

\begin{lstlisting}
array[0] = 100
map["key"] = "value"
\end{lstlisting}

\section{Expression Statements}

Any expression can be used as a statement:

\begin{lstlisting}
io.println("Hello")     // Function call
process(data)           // Function call, ignoring return
x + y                   // Valid but pointless
\end{lstlisting}

\begin{note}
The compiler may warn about expression statements that have no side effects.
\end{note}

\section{If Statements}

\begin{lstlisting}
if condition {
    // executed if condition is true
}

if condition {
    // true branch
} else {
    // false branch
}

if condition1 {
    // first
} else if condition2 {
    // second
} else {
    // default
}
\end{lstlisting}

\subsection{Condition Binding}

\begin{lstlisting}
// Bind and check in one statement
if value = get_optional_value(); value != nil {
    io.println(value)
}
\end{lstlisting}

\section{For Statements}

\subsection{Range-Based For}

\begin{lstlisting}
// Exclusive range
for i in 0..10 {
    io.println(i)  // 0 through 9
}

// Inclusive range
for i in 0..=10 {
    io.println(i)  // 0 through 10
}
\end{lstlisting}

\subsection{Collection Iteration}

\begin{lstlisting}
// Array iteration
for item in items {
    io.println(item)
}

// With index
for i, item in items {
    io.println("${i}: ${item}")
}

// Map iteration
for key, value in scores {
    io.println("${key}: ${value}")
}
\end{lstlisting}

\subsection{String Iteration}

\begin{lstlisting}
for char in "hello" {
    io.println(char)  // Each UTF-8 character
}

for i, char in "hello" {
    io.println("${i}: ${char}")
}
\end{lstlisting}

\section{While Statements}

\begin{lstlisting}
while condition {
    // body
}

// Infinite loop
while true {
    if should_stop() {
        break
    }
}
\end{lstlisting}

\section{Match Statements}

\subsection{Basic Matching}

\begin{lstlisting}
match value {
    1 => io.println("one")
    2 => io.println("two")
    3 => io.println("three")
    _ => io.println("other")
}
\end{lstlisting}

\subsection{Multiple Patterns}

\begin{lstlisting}
match value {
    1 | 2 | 3 => io.println("small")
    4 | 5 | 6 => io.println("medium")
    _ => io.println("large")
}
\end{lstlisting}

\subsection{Range Patterns}

\begin{lstlisting}
match score {
    90..=100 => io.println("A")
    80..90 => io.println("B")
    70..80 => io.println("C")
    60..70 => io.println("D")
    _ => io.println("F")
}
\end{lstlisting}

\subsection{Guard Clauses}

\begin{lstlisting}
match value {
    n if n < 0 => io.println("negative")
    n if n == 0 => io.println("zero")
    n if n > 0 => io.println("positive")
}
\end{lstlisting}

\subsection{Destructuring Patterns}

\begin{lstlisting}
// Struct destructuring
match user {
    User{name: "admin", active: true} => io.println("Active admin")
    User{name: n, active: false} => io.println("Inactive: " + n)
    _ => io.println("Regular user")
}

// Tuple destructuring
match point {
    (0, 0) => io.println("origin")
    (x, 0) => io.println("on x-axis at " + to_string(x))
    (0, y) => io.println("on y-axis at " + to_string(y))
    (x, y) => io.println("at (" + to_string(x) + ", " + to_string(y) + ")")
}

// Enum destructuring
match result {
    Result.Ok(value) => io.println("Got: " + to_string(value))
    Result.Err(err) => io.println("Error: " + err.message)
}
\end{lstlisting}

\subsection{Match Exhaustiveness}

\begin{lstlisting}
enum Direction { North, South, East, West }

match direction {
    Direction.North => go_north()
    Direction.South => go_south()
    Direction.East => go_east()
    Direction.West => go_west()
    // No _ needed: all cases covered
}
\end{lstlisting}

\begin{note}
The compiler verifies that match expressions are exhaustive. If not all cases are covered and there is no wildcard (\code{\_}), the compiler produces an error.
\end{note}

\section{Block Statements}

\begin{lstlisting}
{
    x = 10
    y = 20
    io.println(x + y)
}
\end{lstlisting}

Blocks create a new scope. Variables declared inside are not visible outside.

\section{Return Statements}

\begin{lstlisting}
fn add(a: int, b: int) -> int {
    return a + b
}

fn early_return(x: int) -> int {
    if x < 0 {
        return 0
    }
    return x * 2
}

// Multiple returns
fn divide(a: int, b: int) -> (int, Error?) {
    if b == 0 {
        return 0, Error{message: "division by zero"}
    }
    return a / b, nil
}
\end{lstlisting}

\subsection{Implicit Return}

The last expression in a function is implicitly returned:

\begin{lstlisting}
fn add(a: int, b: int) -> int {
    a + b  // No return keyword needed
}
\end{lstlisting}

\section{Break Statements}

\begin{lstlisting}
for i in 0..100 {
    if i > 10 {
        break
    }
    io.println(i)
}

while true {
    if done() {
        break
    }
}
\end{lstlisting}

\subsection{Labeled Break}

\begin{lstlisting}
outer: for i in 0..10 {
    for j in 0..10 {
        if i * j > 25 {
            break outer
        }
    }
}
\end{lstlisting}

\section{Continue Statements}

\begin{lstlisting}
for i in 0..10 {
    if i % 2 == 0 {
        continue  // Skip even numbers
    }
    io.println(i)
}
\end{lstlisting}

\subsection{Labeled Continue}

\begin{lstlisting}
outer: for i in 0..5 {
    for j in 0..5 {
        if j == 2 {
            continue outer
        }
        io.println("${i}, ${j}")
    }
}
\end{lstlisting}

\section{Import Statements}

Import statements must appear at the top of the file:

\begin{lstlisting}
import "io"
import "json"
import db from "haira/postgres"
import { User, Post } from "models"
\end{lstlisting}

See Chapter~\ref{ch:modules} for complete import semantics.

\section{Defer Statements}

\keyword{defer} schedules a statement to run when the current scope exits:

\begin{lstlisting}
fn read_file(path: string) -> (string, Error?) {
    file, err = fs.open(path)
    if err != nil {
        return "", err
    }
    defer fs.close(file)

    content, err = fs.read_all(file)
    if err != nil {
        return "", err  // file is closed here
    }

    return content, nil  // file is closed here too
}
\end{lstlisting}

\subsection{Multiple Defers}

Multiple \keyword{defer} statements execute in reverse order (LIFO):

\begin{lstlisting}
fn example() {
    defer io.println("first")   // Runs third
    defer io.println("second")  // Runs second
    defer io.println("third")   // Runs first
}
// Output: third, second, first
\end{lstlisting}

\section{Statement Grammar}

\begin{lstlisting}[language=EBNF]
statement = variable_decl
          | assignment
          | expression_stmt
          | if_stmt
          | for_stmt
          | while_stmt
          | match_stmt
          | return_stmt
          | break_stmt
          | continue_stmt
          | defer_stmt
          | block ;

variable_decl = identifier [ ":" type ] "=" expression
              | identifier_list "=" expression_list ;

assignment = lvalue assign_op expression ;

assign_op = "=" | "+=" | "-=" | "*=" | "/=" ;

lvalue = identifier | field_access | index_access ;

expression_stmt = expression ;

if_stmt = "if" [ simple_stmt ";" ] expression block
          [ "else" ( if_stmt | block ) ] ;

for_stmt = "for" [ identifier "," ] identifier "in" expression block ;

while_stmt = "while" expression block ;

match_stmt = "match" expression "{" { match_arm } "}" ;

match_arm = pattern [ "if" expression ] "=>" ( expression | block ) ;

return_stmt = "return" [ expression_list ] ;

break_stmt = "break" [ identifier ] ;

continue_stmt = "continue" [ identifier ] ;

defer_stmt = "defer" statement ;

block = "{" { statement } "}" ;
\end{lstlisting}
